% ============================================================================
% Section 3.4 — Memory-Bounded Verify-Time Refresh
% Insert after Section 3.3, before Section 4 (Experiments).
% Packages: amsmath, algorithm, algpseudocode.
% NOTE on cross-refs: replace \eqref{eq:merge}, \eqref{eq:slice},
%   \eqref{eq:speedup} with the actual labels of your Eq. (4), Eq. (5), Eq. (2).
% Notation reused from 3.1-3.3:
%   n experts {E_i}, L layers, counts c_i over window W, top-M set M,
%   merged experts \hat{E} (K=1) / \hat{E}^{(j)} over slices Q_j (K>1),
%   budget rho = K/n, speedup S with acceptance term tau(gamma) [Eq. (2)].
% ============================================================================
\subsection{Memory-Bounded Verify-Time Refresh}
\label{sec:verify-merge}

Sections~\ref{sec:topm}--\ref{sec:cluster} construct the merged experts
$\widehat{\mathcal{E}}$ and argue that their per-cycle refresh
(Eq.~\eqref{eq:merge}, Eq.~\eqref{eq:slice}) is \emph{interconnect-free}:
verification already fetches the activated experts, so re-evaluating the
count-weighted merge over the top-$M$ constituents $\mathcal{M}$ only reuses
matrices that are already GPU-resident. This argument is sound only if the
constituents are \emph{still} resident at the moment the merge runs. On a
production offloading engine under the tight budget $\rho$, they are not: when
the expert cache is full during batched verification, the engine evicts each
activated expert back to host \emph{immediately after its use}. A na\"ive refresh
deferred to the cycle boundary therefore finds every constituent of
Eq.~\eqref{eq:merge} already evicted and must re-fetch them over PCIe---
reintroducing exactly the migration cost the merged draft was designed to remove.
Making the refresh genuinely free thus requires \emph{co-scheduling} the merge
with verification so that the constituents are resident when the merge consumes
them, while never exceeding the budget. We realise this with a four-stage memory
schedule, $\mathrm{P1}$--$\mathrm{P4}$, applied in dependency order.

\paragraph{P1: Phase-exclusive residency.}
Drafting and verification never run concurrently, so their working sets need not
co-reside. We give each phase exclusive ownership of the expert-memory budget:
at the end of a draft phase we flush the $K$ pinned merged experts
$\{\widehat{\mathcal{E}}\}$, returning their footprint to the verification cache,
and symmetrically reclaim that cache at the start of the next draft. Since the
budget $\rho=K/n$ that pins the merged experts would otherwise permanently shrink
the cache available to verification, phase-exclusive ownership both keeps peak
memory within $\rho$ and, by leaving the verification cache less full, reduces how
often the engine falls back to the immediate post-use eviction that destroys
residency.

\paragraph{P2: Constituent-last ordering with per-layer eviction.}
Within each verification layer we reorder expert fetches so that the
non-constituent experts are dispatched first and the merge constituents
$\mathcal{M}$ last, maximising the interval over which the constituents remain
resident before the merge consumes them. Once a layer has been both verified and
merged, we evict its experts in a single per-layer pass, bounding the live cache
to roughly one layer's working set so the schedule fits the budget at every
depth. This ordering also serialises the merge against the engine's eviction,
guaranteeing that the in-place merge of P3 never reads a partially-evicted
constituent.

\paragraph{P3: In-place verify-time merge.}
With residency guaranteed by P1 and P2, we evaluate the count-weighted merge
(Eq.~\eqref{eq:merge} for $K{=}1$, Eq.~\eqref{eq:slice} for $K{>}1$) in place,
per layer, the moment layer $\ell$'s verification completes, directly from the
resident constituent weights through a single fused on-device kernel. No
host-to-device transfer is incurred. Because the counts $c_i$ that drive the
merge are exactly those produced by the verification pass, relocating the merge
from the cycle boundary into the pass is \emph{output-identical} to the refresh
of Sections~\ref{sec:topm}--\ref{sec:cluster}: it changes \emph{when} the merged
experts are built, not \emph{what} they are, leaving the draft weights---and hence
the acceptance term $\tau(\gamma)$ in Eq.~\eqref{eq:speedup}---unchanged. P2 and
P3 share a single per-layer hook and are realised as one mechanism.

\paragraph{P4: Merge--fetch overlap.}
The in-place merge of P3 is a synchronous on-device computation. As an orthogonal
extension, its latency can be hidden by issuing the merge on a side stream that
overlaps the host-to-device fetch of the following layer's experts, at the cost
of pinning each constituent until its side-stream merge completes. We describe P4
for completeness and leave its implementation and evaluation to future work.%
\footnote{P1--P3 are enabled in all reported results; P4 is a design extension
that we do not evaluate.}

Together, P1--P3 turn the ``free refresh'' of Sections~\ref{sec:topm}--%
\ref{sec:cluster} from an asymptotic argument into a property that holds on a real
offloading engine at the budget $\rho$: every speculative cycle rebuilds all
merged experts with zero re-fetch, so the only expert migration that remains is
the verification traffic the target would incur regardless.

\begin{algorithm}[t]
\caption{One speculative cycle with memory-bounded refresh (P1--P3)}
\label{alg:verify-merge}
\begin{algorithmic}[1]
\State \textbf{Draft:} generate $\gamma$ tokens using resident
       $\{\widehat{\mathcal{E}}\}$ \Comment{no expert migration}
\State flush $\{\widehat{\mathcal{E}}\}$, reclaim verification cache \Comment{P1}
\For{each layer $\ell = 1,\dots,L$ in the verification pass}
  \State evaluate gate; obtain counts $c_i$ and top-$M$ set $\mathcal{M}$
         \Comment{byproduct of verification}
  \State dispatch non-constituent experts, then $\mathcal{M}$ \Comment{P2}
  \State verify layer $\ell$
  \State build $\widehat{\mathcal{E}}_\ell$ from resident weights via
         Eq.~\eqref{eq:merge}/\eqref{eq:slice} \Comment{P3, 0 re-fetch}
  \State evict layer $\ell$'s experts \Comment{P2}
\EndFor
\State pin $\{\widehat{\mathcal{E}}\}$ for the next draft \Comment{P1}
\end{algorithmic}
\end{algorithm}
